% Physical connection change: one split event, two local branches, one output.
% The event sketch keeps only the before/after cube pair; the left branch
% updates local connections and surfaces, the right branch keeps the material
% motion history and the surface-existence history apart.  Visibility gaps are
% not drawn as surface birth or as a physical disconnection.
\begingroup
% Shared editable vector language for the five Section 3.2 technical routes.
%
% The figures are explanatory diagrams, not experimental results.  A single
% visual grammar is used throughout:
%   * a dark header band names a block; white type is reserved for headers;
%   * a numbered marker gives the reading order of a route, and the marker of a
%     stage that carries the study's own mechanism is filled;
%   * white cards with a thin theme-coloured rule hold stage titles and bodies;
%   * tinted panels hold mechanism sketches, and a heavier frame marks the
%     representation or relation that is handed to the next study;
%   * solid arrows are information flow, dashed arrows are observation
%     feedback or a model-predicted interval.
% Colour is always paired with a marker, line style or label so that the
% grayscale proof remains readable.
\linespread{1}\selectfont
\usetikzlibrary{arrows.meta,calc,positioning,matrix,decorations.pathreplacing}

% ---- palette -------------------------------------------------------------
\definecolor{trBlueDark}{HTML}{1F4E79}
\definecolor{trBlueMid}{HTML}{2F6690}
\definecolor{trBlueSoft}{HTML}{7FA6C4}
\definecolor{trBluePale}{HTML}{E8F1F8}
\definecolor{trBlueWash}{HTML}{F4F8FB}
\definecolor{trGreenDark}{HTML}{2E5E40}
\definecolor{trGreenMid}{HTML}{3E7A52}
\definecolor{trGreenSoft}{HTML}{8FB39A}
\definecolor{trGreenPale}{HTML}{E9F2EB}
\definecolor{trGreenWash}{HTML}{F4F9F5}
\definecolor{trOchreDark}{HTML}{7A5018}
\definecolor{trOchreMid}{HTML}{9E6B27}
\definecolor{trOchreSoft}{HTML}{C9A46B}
\definecolor{trOchrePale}{HTML}{F8EEDC}
\definecolor{trOchreWash}{HTML}{FCF7EE}
\definecolor{trVioletDark}{HTML}{4A3A6A}
\definecolor{trVioletMid}{HTML}{67528E}
\definecolor{trVioletSoft}{HTML}{A294BE}
\definecolor{trVioletPale}{HTML}{EFEAF6}
\definecolor{trVioletWash}{HTML}{F8F5FC}
\definecolor{trInkSoft}{HTML}{4D4D4D}
\definecolor{trRule}{HTML}{9AA5AE}

\newcommand{\trTheme}[1]{%
  \colorlet{trDark}{tr#1Dark}%
  \colorlet{trMid}{tr#1Mid}%
  \colorlet{trSoft}{tr#1Soft}%
  \colorlet{trPale}{tr#1Pale}%
  \colorlet{trFill}{tr#1Pale}%
  \colorlet{trWash}{tr#1Wash}%
}
\trTheme{Blue}

% ---- shared type scale for every Part 3 figure ----------------------------
% Set sizes here rather than overriding individual diagrams.  The report uses
% the natural figure width, so these remain the printed sizes; subscripts use
% the normal mathematical script sizes.
\newcommand{\trFontHead}{\songti\bfseries\fontsize{9.2pt}{11.2pt}\selectfont}
\newcommand{\trFontSubhead}{\songti\bfseries\fontsize{8.2pt}{10.2pt}\selectfont}
\newcommand{\trFontBody}{\songti\fontsize{8.2pt}{10.2pt}\selectfont}
\newcommand{\trFontSmall}{\songti\fontsize{8pt}{10pt}\selectfont}
\newcommand{\trFontTiny}{\songti\fontsize{8pt}{10pt}\selectfont}
\newcommand{\trFontMath}{\fontsize{9.2pt}{11.2pt}\selectfont}
\newcommand{\trFontMathSmall}{\fontsize{8.5pt}{10.4pt}\selectfont}
\newcommand{\trFontMarker}{\songti\bfseries\fontsize{8pt}{10pt}\selectfont}
\newcommand{\trFontSubcaption}{\songti\mdseries\fontsize{9.2pt}{11.2pt}\selectfont}

% ---- base styles ---------------------------------------------------------
\tikzset{
  trBase/.style={x=1cm,y=1cm,text=black,line cap=round,line join=round,
    font=\trFontBody,
    every node/.style={align=center,outer sep=0pt,inner sep=0pt}},
  trTitle/.style={font=\trFontHead},
  trTitleSmall/.style={font=\trFontSubhead},
  trHead/.style={font=\trFontHead,text=white},
  trBody/.style={font=\trFontBody},
  trNote/.style={font=\trFontSmall},
  trTiny/.style={font=\trFontTiny},
  trLabel/.style={font=\trFontSmall},
  trMath/.style={font=\trFontMath},
  trMathSmall/.style={font=\trFontMathSmall},
  trNum/.style={circle,draw=trDark,fill=trDark,text=white,
    line width=0.5pt,inner sep=0pt,minimum size=0.36cm,
    font=\trFontMarker},
  trNumOpen/.style={trNum,fill=white,text=trDark},
  trCard/.style={fill=white,draw=trMid,line width=0.55pt,rounded corners=1.2pt},
  trFrame/.style={draw=black,line width=0.65pt},
  trArrow/.style={draw=trDark,line width=0.8pt,
    -{Latex[length=1.7mm,width=1.15mm]}},
  trArrowThin/.style={draw=trDark,line width=0.6pt,
    -{Latex[length=1.4mm,width=0.95mm]}},
  trFlow/.style={draw=black!72,line width=0.7pt,
    -{Latex[length=1.6mm,width=1.05mm]}},
  trFeedback/.style={draw=trDark,line width=0.8pt,dash pattern=on 3pt off 2pt,
    -{Latex[length=1.7mm,width=1.15mm]}},
  trPredict/.style={draw=trMid,line width=0.75pt,dash pattern=on 2.2pt off 1.6pt,
    -{Latex[length=1.6mm,width=1.05mm]}},
  trGuide/.style={draw=black!35,line width=0.45pt},
  trTag/.style={draw=trMid,fill=white,line width=0.45pt,inner sep=2.2pt,
    rounded corners=0.8pt,font=\trFontSmall},
  trAccent/.style={draw=trMid,line width=0.9pt},
  trAccentThin/.style={draw=trMid,line width=0.6pt},
  trPossible/.style={draw=trSoft,line width=0.7pt,dash pattern=on 1.8pt off 1.4pt},
  trTimeLine/.style={draw=black!55,line width=0.6pt,
    -{Latex[length=1.5mm,width=0.95mm]}},
  trDown/.style={draw=trDark,line width=0.8pt,
    -{Latex[length=1.6mm,width=1.05mm]}},
  trDouble/.style={draw=trDark,line width=0.8pt,
    {Latex[length=1.5mm,width=1.0mm]}-{Latex[length=1.5mm,width=1.0mm]}},
  trDoubleSmall/.style={draw=trDark,line width=0.7pt,
    {Latex[length=1.1mm,width=0.75mm]}-{Latex[length=1.1mm,width=0.75mm]}}
}

% ---- reusable pieces -----------------------------------------------------
% Header band.  #1 left, #2 right, #3 bottom, #4 top, #5 text.
\newcommand{\trHeader}[5]{%
  \path[fill=trDark] (#1,#3) rectangle (#2,#4);
  \node[trHead] at ({#1/2+#2/2},{#3/2+#4/2}) {#5}
}

% ---- compact detailed-route primitives (2026-10-30) ----------------------
% The four detailed routes use a downward y axis so that a route reads from
% the input strip at the top to the output strip at the bottom.  Their canvases
% are 15.56 cm wide and only as tall as the content; the styles below keep the
% four redraws on the same type scale as the overview (Fig. 11).
\tikzset{
  trRouteBase/.style={x=1cm,y=-1cm,text=black,line cap=round,line join=round,
    font=\trFontBody,
    every node/.style={align=center,outer sep=0pt,inner sep=0pt}},
  trRouteHead/.style={font=\trFontHead,text=white},
  trRouteBody/.style={font=\trFontBody},
  trRouteSmall/.style={font=\trFontSmall},
  trRouteTiny/.style={font=\trFontTiny},
  trRouteMath/.style={font=\trFontMath},
  trRouteMathSmall/.style={font=\trFontMathSmall},
  trRouteFlow/.style={draw=trDark,line width=0.72pt,
    -{Latex[length=1.5mm,width=1.0mm]}},
  trRouteThin/.style={draw=trDark,line width=0.6pt,
    -{Latex[length=1.3mm,width=0.9mm]}},
  trRouteBack/.style={draw=trDark,line width=0.65pt,
    dash pattern=on 2.6pt off 1.8pt,-{Latex[length=1.4mm,width=0.95mm]}},
  trRouteRule/.style={draw=trMid,line width=0.55pt},
  trRouteMotion/.style={draw=trBlueDark,line width=0.95pt},
  trRoutePossible/.style={draw=black!52,line width=0.65pt,
    dash pattern=on 1.8pt off 1.4pt},
  trRouteCard/.style={fill=white,draw=trMid,line width=0.6pt},
  trRouteStrip/.style={fill=trWash,draw=trMid,line width=0.55pt},
  trRouteFocus/.style={fill=trPale,draw=trDark,line width=0.7pt}
}

\trTheme{Ochre}
% ---- local panel helpers (scoped to this figure) -------------------------
\newcommand{\trkPanel}[5]{%
  \path[fill=white,draw=trMid,line width=0.6pt] (#1,#2) rectangle (#3,#4);
  \path[fill=trDark] (#1,#2) rectangle (#3,{#2+0.56});
  \node[trRouteHead] at ({#1/2+#3/2},{#2+0.28}) {#5};}
\newcommand{\trkFocus}[5]{%
  \path[fill=trPale,draw=trDark,line width=0.7pt] (#1,#2) rectangle (#3,#4);
  \path[fill=trDark] (#1,#2) rectangle (#3,{#2+0.56});
  \node[trRouteHead] at ({#1/2+#3/2},{#2+0.28}) {#5};}
\newcommand{\trkPanelMain}[5]{%
  \path[fill=white,draw=trMid,line width=0.6pt] (#1,#2) rectangle (#3,#4);
  \path[fill=trDark] (#1,#2) rectangle (#3,{#2+0.56});
  \node[trRouteHead] at ({#1/2+#3/2},{#2+0.28}) {#5};}
% ---- local block primitives ----------------------------------------------
% front face + top face; the right side is added only where it is visible
\newcommand{\trkBody}[6]{%
  \path[fill=trBluePale,draw=trBlueDark,line width=0.6pt]
    (#1,#2) rectangle ({#1+#3},{#2+#4});
  \path[fill=trBlueWash,draw=trBlueDark,line width=0.6pt]
    (#1,#2) -- ({#1+#5},{#2+#6}) -- ({#1+#3+#5},{#2+#6}) --
    ({#1+#3},#2) -- cycle;}
\newcommand{\trkBodySide}[6]{%
  \trkBody{#1}{#2}{#3}{#4}{#5}{#6}%
  \path[fill=trBluePale!70,draw=trBlueDark,line width=0.6pt]
    ({#1+#3},#2) -- ({#1+#3+#5},{#2+#6}) --
    ({#1+#3+#5},{#2+#4+#6}) -- ({#1+#3},{#2+#4}) -- cycle;}
% new cut face: a vertical face at the split, drawn as a parallelogram
\newcommand{\trkCut}[5]{%
  \path[fill=trOchrePale,draw=trOchreDark,line width=0.85pt]
    (#1,#2) -- ({#1+#4},{#2+#5}) -- ({#1+#4},{#2+#3+#5}) --
    (#1,{#2+#3}) -- cycle;}
\begin{tikzpicture}[trRouteBase]
  \path[use as bounding box] (0.12,0) rectangle (15.68,8.27);

  % ======================= event: one split ==============================
  \trkPanelMain{0.15}{0.15}{15.65}{3.00}{切分事件与连接判断}

  % intact body and the connection that survives in the stable region
  \trkBodySide{1.30}{1.38}{1.15}{1.10}{0.28}{-0.18}
  \draw[trOchreDark,line width=1.5pt] (1.42,1.68) -- (2.33,1.68);
  \draw[trOchreDark,line width=0.8pt] (1.42,1.60) -- (1.42,1.76);
  \draw[trOchreDark,line width=0.8pt] (2.33,1.60) -- (2.33,1.76);
  \node[trRouteSmall] at (1.95,1.00) {稳定区连接};

  % the one physical event
  \draw[trRouteFlow] (3.05,1.83) -- (4.38,1.83);
  \node[trRouteSmall] at (3.71,1.54) {切分};

  % the two halves: same height and depth, complementary cut faces
  \trkBody{4.55}{1.38}{0.55}{1.10}{0.20}{-0.15}
  \trkCut{5.10}{1.38}{1.10}{0.20}{-0.15}
  \trkBodySide{5.72}{1.38}{0.55}{1.10}{0.20}{-0.15}
  \trkCut{5.72}{1.38}{1.10}{-0.20}{-0.15}

  % the released connection across the split
  \draw[trRoutePossible,draw=black!60] (5.30,1.74) -- (5.52,1.74);
  \draw[black!65,line width=1.0pt] (5.33,1.66) -- (5.49,1.82)
    (5.33,1.82) -- (5.49,1.66);
  \node[trRouteSmall] at (5.41,1.00) {失效连接};

  % the cut faces are the only new surfaces
  \draw[trOchreDark,line width=0.5pt] (5.30,2.56) -- (5.30,2.36);
  \draw[trOchreDark,line width=0.5pt] (5.52,2.56) -- (5.52,2.36);
  \node[trRouteSmall,text=trOchreDark] at (5.40,2.78)
    {新生表面};

  % evidence and the visibility boundary, kept to the right of the sketch
  \path[trRouteStrip] (7.85,1.08) rectangle (15.50,2.62);
  \fill[trOchreDark] (7.98,1.395) rectangle ++(0.09,0.09);
  \node[trRouteSmall,anchor=west] at (8.16,1.44)
    {判据：分离运动 · 轮廓／深度 · 事件};
  \draw[black!60,fill=white,line width=0.55pt]
    (7.98,1.805) rectangle ++(0.09,0.09);
  \node[trRouteSmall,anchor=west] at (8.16,1.85)
    {边界：观测中断不判新生/断连};
  \draw[black!60,fill=white,line width=0.55pt] (8.025,2.26) circle (0.055);
  \node[trRouteSmall,anchor=west] at (8.16,2.26)
    {对象：同一物体};

  % ================= split to the two local branches =====================
  \draw[draw=trDark,line width=0.72pt] (7.90,3.00) -- (7.90,3.20) --
    (3.85,3.20) -- (11.95,3.20);
  \draw[trRouteFlow] (3.85,3.20) -- (3.85,3.36);
  \draw[trRouteFlow] (11.95,3.20) -- (11.95,3.36);

  % ================= left branch: local update ===========================
  \trkPanel{0.15}{3.40}{7.55}{6.60}{局部连接与表面更新}
  \draw[trOchreDark,line width=1.4pt] (0.46,4.24) -- (0.70,4.24);
  \node[trRouteBody,anchor=west] at (0.85,4.24)
    {保持稳定区关系};
  \draw[black!65,line width=0.9pt] (0.525,4.505) -- (0.635,4.615)
    (0.525,4.615) -- (0.635,4.505);
  \node[trRouteBody,anchor=west] at (0.85,4.56)
    {解除失效连接};
  \fill[trOchreDark] (0.53,4.83) rectangle ++(0.10,0.10);
  \node[trRouteBody,anchor=west] at (0.85,4.88)
    {更新表面与新生边界};

  \path[fill=trOchrePale,draw=trOchreSoft,line width=0.55pt]
    (1.65,5.16) rectangle (6.05,5.64);
  \node[trRouteMath] at (3.85,5.40) {$g_{ij}e_{ij}(t)$};
  \node[trRouteSmall] at (3.85,5.94)
    {$g_{ij}$：共享权重，非物理连接};
  \node[trRouteSmall] at (3.85,6.33)
    {$e_{ij}(t)$：$\mathcal E(t)$ 物理连接支持度};

  % ================= right branch: two histories =========================
  \trkPanel{8.25}{3.40}{15.65}{6.60}{材料运动与表面历史}
  \node[trRouteSmall,text=trOchreDark] at (12.525,4.20) {形成区间待定};
  \node[trRouteSmall,text=trBlueDark] at (9.35,4.87)
    {有效关联下延续};

  % formation-time interval: the surface history is pending inside it
  \path[fill=trOchrePale] (12.10,4.40) rectangle (12.95,5.85);
  \draw[trOchreSoft,line width=0.55pt,dash pattern=on 2pt off 1.5pt]
    (12.10,4.40) -- (12.10,5.85) (12.95,4.40) -- (12.95,5.85);

  % material motion history: solid where association support is available
  \draw[draw=trBlueDark,line width=0.95pt,
    -{Latex[length=1.3mm,width=0.9mm]}] plot[smooth] coordinates
    {(10.55,4.60) (11.20,4.50) (12.10,4.54) (12.95,4.49)
     (13.80,4.56) (14.50,4.48) (15.30,4.55)};
  \foreach \xx/\yy in {10.62/4.59,11.20/4.50,11.70/4.53,12.30/4.52,13.15/4.49,13.80/4.54,14.50/4.52,15.00/4.49}{
    \fill[trBlueDark] (\xx,\yy) circle (0.045);}
  \node[trRouteSmall,text=trBlueDark] at (9.35,4.42) {材料运动历史};

  % new-surface history: pending (hollow) before the confirmed formation
  \draw[draw=trOchreDark,line width=0.85pt,dash pattern=on 2pt off 1.5pt]
    (12.10,5.42) -- (12.95,5.42);
  \foreach \xx in {12.10,12.38,12.66}{
    \draw[trOchreDark,fill=white,line width=0.6pt]
      (\xx,5.42) circle (0.045);}
  \draw[draw=trOchreDark,line width=0.85pt,
    -{Latex[length=1.3mm,width=0.9mm]}] (12.95,5.42) -- (15.30,5.42);
  \foreach \xx in {12.95,13.60,14.25,14.90}{
    \fill[trOchreDark] (\xx,5.42) circle (0.045);}
  \node[trRouteSmall,text=trOchreDark] at (9.35,5.38) {新生表面历史};
  \node[trRouteMathSmall,text=trOchreDark] at (9.35,5.78) {$I_i^{\mathrm{surf}}$};
  \node[trRouteSmall,text=trOchreDark] at (13.85,5.14) {形成后记录};

  % shared time axis
  \draw[draw=black!60,line width=0.55pt,
    -{Latex[length=1.2mm,width=0.8mm]}] (10.55,6.25) -- (15.320,6.25);
  \node[trRouteSmall,text=black!75] at (15.386,6.25) {$t$};

  % ================= merge and output ====================================
  \draw[draw=trDark,line width=0.72pt] (3.85,6.60) -- (3.85,6.78) --
    (11.95,6.78) -- (11.95,6.60);
  \draw[trRouteFlow] (7.90,6.78) -- (7.90,6.96);
  \trkFocus{2.55}{7.00}{13.25}{8.12}
    {更新几何 $\mathcal G(t)$ 与物理连接 $\mathcal E(t)$}
  \node[trRouteSmall,text=black!75] at (7.90,7.84)
    {辅助：双历史与待定连接};
% \ignorespaces absorbs the wrapper's newline after \input; without it the
% standalone minipage adds an empty bottom line to the preview/export.
\end{tikzpicture}\ignorespaces
\endgroup
